% !TeX root = tcc.tex
% !TeX encoding = UTF-8
% !BIB TS-program = biber

% ─────────────────────────────────────────────────────────────
%  MIDI HERO — Monografia de Graduação
%  Samuel Moreira da Silva — DECAT/UFOP — 2026
% ─────────────────────────────────────────────────────────────

\documentclass[
  12pt,
  openright,
  twoside,
  a4paper,
  english,
  brazil
]{abntex2}

% ── Pacotes básicos ──────────────────────────────────────────
\usepackage{graphicx}
\usepackage[table,xcdraw]{xcolor}
\usepackage{booktabs}
\graphicspath{{figuras/}}
\DeclareGraphicsExtensions{.pdf,.eps,.svg,.png,.jpg,.bmp}

\usepackage{microtype}
\usepackage{amsmath,amssymb,unicode-math}
\usepackage{relsize}
\usepackage{verbatim}
\usepackage{pdfpages}
\usepackage{csquotes}
\usepackage{supertabular}
\usepackage{multicol}
\usepackage{multirow}
\usepackage{lipsum}

% ── Código-fonte ─────────────────────────────────────────────
\usepackage{listings}
\lstset{
  basicstyle=\ttfamily\small,
  breaklines=true,
  frame=single,
  language=Python,
  keywordstyle=\bfseries,
  commentstyle=\itshape\color{gray},
  showstringspaces=false
}

% ── Citações (autor-data ABNT) ────────────────────────────────
\usepackage[
  backend=biber,
  style=abnt,
  sccite,
  ittitles,
  citecount,
  scbib,
  justify,
  noslsn,
  repeatfields,
  sorting=nty
]{biblatex}

\addbibresource{referencias.bib}

\DeclareFieldFormat{url}{\bibstring{urlfrom}\addcolon\addspace\url{#1}}
\DeclareFieldFormat{urldate}{\bibstring{urlseen}\addcolon\addspace #1}

\makeatletter
\@addtoreset{footnote}{chapter}
\makeatother

% ── Metadados ────────────────────────────────────────────────
\titulo{MIDI HERO: SISTEMA GAMIFICADO PARA PRÁTICA DE PIANO E TECLADO COM SUPORTE AO PROTOCOLO MIDI}
\autor{SAMUEL MOREIRA DA SILVA}
\local{Ouro Preto}
\data{\the\year}
\orientador{Profa. Dra. Adrielle de Carvalho Santana}
\instituicao{Universidade Federal de Ouro Preto}
\tipotrabalho{Monografia de Graduação}
\preambulo{%
  Trabalho apresentado ao Colegiado do Curso de Engenharia de
  Controle e Automação da Universidade Federal de Ouro Preto como
  parte dos requisitos para a obtenção do Grau de Engenheiro(a)
  de Controle e Automação.}

% ── PDF ──────────────────────────────────────────────────────
\definecolor{blue}{RGB}{41,5,195}
\makeatletter
\hypersetup{
  pdftitle={\@title},
  pdfauthor={\@author},
  pdfsubject={\imprimirpreambulo},
  pdfcreator={LaTeX with abnTeX2},
  pdfkeywords={ufop}{midi}{piano}{gamificação}{pygame}{python}{decat},
  colorlinks=true,
  linkcolor=blue,
  citecolor=blue,
  filecolor=magenta,
  urlcolor=blue
}
\makeatother

% ── Espaçamento ──────────────────────────────────────────────
\setlength{\parindent}{1.3cm}
\setlength{\parskip}{0.2cm}

\makeindex

% ══════════════════════════════════════════════════════════════
\begin{document}
\frenchspacing

% ─────────────────────────────────────────────────────────────
%  CAPA
% ─────────────────────────────────────────────────────────────
\begin{capa}
  \thispagestyle{empty}
  \centering
  \begin{minipage}{1\linewidth}
    \centering
    \begin{tabular}{ccc}
      \begin{tabular}{c}\\
        \includegraphics[width=0.9cm]{logo-universidade.jpg}
      \end{tabular}
      &
      \begin{tabular}{c}\\
        {\large \imprimirinstituicao}\\
        {\large Escola de Minas}\\
        {\large CECAU --- Colegiado do Curso de}\\
        {\large Engenharia de Controle e Automação}
      \end{tabular}
      &
      \begin{tabular}{c}\\
        \includegraphics[width=2.1cm]{logo-unidade-2.jpg}
      \end{tabular}
    \end{tabular}
  \end{minipage}

  \vspace*{1cm}
  {\ABNTEXchapterfont\large\imprimirautor}
  \vspace*{\fill}
  {\ABNTEXchapterfont\bfseries\large\imprimirtitulo}
  \vspace*{\fill}
  {\large\imprimirtipotrabalho}
  \vspace*{\fill}
  {\large\imprimirlocal}, {\large\imprimirdata}
  \vspace*{1cm}
\end{capa}

% ─────────────────────────────────────────────────────────────
%  FOLHA DE ROSTO
% ─────────────────────────────────────────────────────────────
\imprimirfolhaderosto*

% ─────────────────────────────────────────────────────────────
%  FICHA CATALOGRÁFICA
% ─────────────────────────────────────────────────────────────
\begin{fichacatalografica}
  \vspace*{\fill}
  \hrule
  \begin{center}
  \begin{minipage}[c]{12.5cm}
    \imprimirautor

    \hspace{0.5cm}\imprimirtitulo / \imprimirautor. --
    \imprimirlocal, \imprimirdata-

    \hspace{0.5cm} p. : il. (algumas color.) ; 30 cm.\\

    \hspace{0.5cm}\imprimirorientadorRotulo~\imprimirorientador\\

    \hspace{0.5cm}
    \parbox[t]{\textwidth}{%
      \imprimirtipotrabalho~--~\imprimirinstituicao, \imprimirdata.}\\

    \hspace{0.5cm}
      1. Gamificação.
      2. Protocolo MIDI.
      3. Educação musical.
      4. Python.
      5. Pygame.
      I. Santana, Adrielle de Carvalho.
      II. Universidade Federal de Ouro Preto.
      III. Escola de Minas.
      IV. Título\\

    \hspace{8.75cm} CDU 004.4:78\\
  \end{minipage}
  \end{center}
  \hrule
\end{fichacatalografica}

% ─────────────────────────────────────────────────────────────
%  FOLHA DE APROVAÇÃO
% ─────────────────────────────────────────────────────────────
\includepdf[scale=0.7,pages=1,nup=1x1,frame]{folha-de-aprovacao}

% ─────────────────────────────────────────────────────────────
%  DEDICATÓRIA
% ─────────────────────────────────────────────────────────────
\begin{dedicatoria}
  \vspace*{\fill}
  \flushright
  \noindent
  \textit{%
    À minha esposa, pelo suporte incondicional,\\
    e aos meus pais, por tudo que me tornei.}
  \vspace*{\fill}
\end{dedicatoria}

% ─────────────────────────────────────────────────────────────
%  AGRADECIMENTOS
% ─────────────────────────────────────────────────────────────
\begin{agradecimentos}
\noindent
Agradeço a Deus, fonte de força e discernimento ao longo desta jornada.

Aos meus pais, pelo incentivo constante e pelos sacrifícios que tornaram
possível a minha formação.

À minha esposa, pela paciência, pelo companheirismo e por não reclamar
das noites em que o computador ficou ligado até tarde.

À Professora Dra. Adrielle de Carvalho Santana, pela orientação, pelas
discussões que moldaram este trabalho e pela disponibilidade mesmo nos
momentos mais críticos do semestre.

Aos colegas do curso de Engenharia de Controle e Automação da UFOP, pela
troca de experiências e pela camaradagem que tornaram os anos de
graduação memoráveis.

Por fim, a todos que, de alguma forma, contribuíram para a realização
deste trabalho. O aprendizado nunca é solitário.
\end{agradecimentos}

% ─────────────────────────────────────────────────────────────
%  EPÍGRAFE
% ─────────────────────────────────────────────────────────────
\begin{epigrafe}
  \vspace*{\fill}
  \begin{flushright}
    \textit{``Pode ser fácil se você ver o mundo de outro jeito.''}\\
    (Capital Inicial)
  \end{flushright}
\end{epigrafe}

% ─────────────────────────────────────────────────────────────
%  RESUMO (português)
% ─────────────────────────────────────────────────────────────
\setlength{\absparsep}{18pt}
\begin{resumo}
\noindent
O aprendizado de piano é uma atividade reconhecida pelos seus benefícios
cognitivos e motores, porém frequentemente inacessível em razão do alto
custo de instrumentos e aulas particulares. Este trabalho apresenta o
desenvolvimento do \textbf{MIDI Hero}, um sistema computacional
gamificado para suporte à prática de piano e teclado que utiliza o
protocolo \textit{MIDI} (\textit{Musical Instrument Digital Interface})
como meio de leitura de partituras e captura da execução do usuário em
tempo real. O sistema exibe barras de notas descendentes sincronizadas à
partitura, avalia a precisão temporal e a duração das notas pressionadas,
e retorna pontuação progressiva e feedback auditivo imediato via síntese
sonora por SoundFont. A implementação foi realizada em Python com a
biblioteca Pygame, organizada em arquitetura modular, e é compatível com
os sistemas operacionais Linux, Windows e macOS. % TODO: completar com resultados da avaliação com usuários.

\vspace{\onelineskip}
\noindent
\textbf{Palavras-chave}: Gamificação. Protocolo MIDI. Educação musical.
Piano. Teclado. Python. Pygame. Feedback visual.
\end{resumo}

% ─────────────────────────────────────────────────────────────
%  ABSTRACT (inglês)
% ─────────────────────────────────────────────────────────────
\begin{resumo}[Abstract]
\begin{otherlanguage*}{english}
\noindent
Learning to play the piano is widely recognized for its cognitive and
motor benefits, yet it remains inaccessible to many due to the high cost
of instruments and private lessons. This work presents the development of
\textbf{MIDI Hero}, a gamified computer system designed to support piano
and keyboard practice using the \textit{MIDI} (\textit{Musical Instrument
Digital Interface}) protocol for sheet-music reading and real-time user
input capture. The system displays descending note bars synchronized with
the score, evaluates timing accuracy and note duration, and provides
progressive scoring alongside immediate auditory feedback via SoundFont
synthesis. The implementation was carried out in Python using the Pygame
library, organized in a modular architecture, and is compatible with
Linux, Windows, and macOS. % TODO: add evaluation results.

\vspace{\onelineskip}
\noindent
\textbf{Keywords}: Gamification. MIDI protocol. Music education.
Piano. Keyboard. Python. Pygame. Visual feedback.
\end{otherlanguage*}
\end{resumo}

% ─────────────────────────────────────────────────────────────
%  LISTAS
% ─────────────────────────────────────────────────────────────
\pdfbookmark[0]{\listfigurename}{lof}
\listoffigures*
\cleardoublepage

\pdfbookmark[0]{\listtablename}{lot}
\listoftables*
\cleardoublepage

\begin{siglas}
  \item[ABNT]  Associação Brasileira de Normas Técnicas
  \item[API]   \textit{Application Programming Interface}
  \item[CLI]   \textit{Command-Line Interface}
  \item[GUI]   \textit{Graphical User Interface}
  \item[HUD]   \textit{Heads-Up Display}
  \item[JSON]  \textit{JavaScript Object Notation}
  \item[MIDI]  \textit{Musical Instrument Digital Interface}
  \item[OGG]   Formato de áudio aberto Ogg Vorbis
  \item[RA]    Realidade Aumentada
  \item[RV]    Realidade Virtual
  \item[SDL]   \textit{Simple DirectMedia Layer}
  \item[SMF]   \textit{Standard MIDI File}
  \item[SUS]   \textit{System Usability Scale}
  \item[UFOP]  Universidade Federal de Ouro Preto
  \item[WAV]   \textit{Waveform Audio File Format}
\end{siglas}

% ─────────────────────────────────────────────────────────────
%  SUMÁRIO
% ─────────────────────────────────────────────────────────────
\pdfbookmark[0]{\contentsname}{toc}
\tableofcontents*
\cleardoublepage

\makeatletter
\@addtoreset{footnote}{chapter}
\makeatother

% ══════════════════════════════════════════════════════════════
%  ELEMENTOS TEXTUAIS
% ══════════════════════════════════════════════════════════════
\textual

% ─────────────────────────────────────────────────────────────
%  CAPÍTULO 1 — INTRODUÇÃO
% ─────────────────────────────────────────────────────────────
\chapter{Introdução}

A prática de instrumentos musicais, em especial o piano, é reconhecida na
literatura científica como uma atividade que promove benefícios cognitivos,
motores e socioemocionais. Pesquisas indicam que o aprendizado musical
favorece o desenvolvimento da coordenação motora fina, da memória de
trabalho, da capacidade de atenção sustentada e da inteligência emocional
\cite{hallam2015,miendlarzewska2014}. No entanto, o acesso ao ensino
formal de piano permanece restrito, sobretudo em países em
desenvolvimento, em decorrência do alto custo de instrumentos e de aulas
particulares.

Paralelamente, a \textbf{gamificação}---entendida como a aplicação de
elementos e princípios de design de jogos em contextos não
lúdicos~\cite{deterding2011}---tem demonstrado potencial para aumentar o
engajamento e a motivação em processos educacionais. Jogos como
\textit{Guitar Hero\textregistered} popularizaram o conceito de
aprendizado musical mediado por interfaces gamificadas, comprovando a
eficácia desse modelo para engajar usuários sem experiência prévia com
instrumentos. A existência de dispositivos \textit{MIDI} de baixo custo
(teclados USB com 49 a 88 teclas), associada à disponibilidade gratuita de
partituras digitais em formato MIDI geradas por \textit{softwares} como o
\textit{MuseScore}~\cite{musescore}, abre uma oportunidade concreta para o
desenvolvimento de sistemas educacionais acessíveis voltados ao
aprendizado de piano.

A motivação pessoal do autor também contribui para a relevância deste
trabalho: cantor em corais por vários anos, Samuel Moreira da Silva sempre
nutriu interesse pelo aprendizado do piano, instrumento ao qual nunca teve
acesso formal. O desenvolvimento deste sistema representa, portanto, tanto
uma contribuição acadêmica quanto uma realização pessoal.

\section{Justificativas e Relevância}

Três eixos principais justificam a realização deste trabalho:

\textbf{(i) Democratização do ensino musical.}
O custo médio de um piano acústico de entrada no Brasil supera
R\,\$\,5.000,00, e aulas particulares custam entre R\,\$\,60,00 e
R\,\$\,200,00 por hora. Teclados MIDI USB, por outro lado, estão
disponíveis por valores a partir de R\,\$\,300,00. Um sistema de software
gratuito e de código aberto, combinado a esses dispositivos de baixo
custo, reduz significativamente a barreira financeira ao aprendizado
musical.

\textbf{(ii) Evidências de eficácia da gamificação.}
\textcite{hamari2014}, em revisão sistemática da literatura, identificaram
que a gamificação produz efeitos positivos na motivação e no engajamento
em ambientes educacionais em 66\% dos estudos analisados. O uso de
feedback visual imediato, progressão por pontuação e metas claras---todos
presentes no sistema desenvolvido---são os elementos mais frequentemente
associados a esses resultados.

\textbf{(iii) Ausência de soluções abertas equivalentes.}
As principais ferramentas existentes para o aprendizado gamificado de
piano (Synthesia, PianoVision\textregistered, Simply Piano, Yousician)
são proprietárias, total ou parcialmente pagas, e não permitem
personalização ou extensão. O sistema proposto é integralmente de código
aberto, permitindo que pesquisadores, educadores e desenvolvedores o
adaptem livremente.

\section{Objetivos}

\subsection{Objetivo Geral}

Desenvolver um sistema computacional gamificado, denominado
\textbf{MIDI Hero}, para suporte à prática de piano e teclado, capaz de
ler arquivos no formato MIDI, gerar feedback visual sincronizado em tempo
real e avaliar a execução do usuário por meio de um sistema de pontuação
progressiva.

\subsection{Objetivos Específicos}

\begin{itemize}
  \item Implementar leitura e análise (\textit{parsing}) de arquivos no
        formato MIDI (\textit{Standard MIDI File});
  \item Desenvolver interface gráfica gamificada com barras de notas
        descendentes sincronizadas à partitura;
  \item Implementar sistema de pontuação progressiva baseado em precisão
        temporal e duração da nota;
  \item Implementar síntese sonora em tempo real para retorno auditivo
        ao usuário durante a prática;
  \item Garantir compatibilidade multiplataforma com Linux, Windows e
        macOS;
  \item Estruturar o sistema em arquitetura modular extensível, de modo
        a viabilizar futuras versões com Realidade Aumentada e suporte a
        dispositivos Android.
\end{itemize}

\section{Metodologia}

O desenvolvimento do trabalho seguiu uma abordagem iterativa e
incremental baseada em protótipos sucessivos, com controle de versão pelo
sistema \textit{Git}~\cite{git2005} e repositório público hospedado no
\textit{GitHub}.

\textbf{Evolução da plataforma de desenvolvimento.}
O projeto foi originalmente concebido para a plataforma
\textit{Unity}\footnote{\url{https://unity.com}} com implementação em
\textit{C\#} e foco em recursos de Realidade Aumentada para dispositivos
\textit{Android}. Ao longo do processo de desenvolvimento, identificou-se
que a curva de aprendizado associada ao \textit{Unity} e ao \textit{C\#}
representaria um obstáculo ao progresso dentro do prazo acadêmico
disponível. Optou-se, então, por uma reformulação estratégica: manter a
visão original de criar um sistema interativo para prática de piano, mas
implementá-lo como um protótipo funcional em Python com a biblioteca
Pygame. Essa decisão permitiu avançar com maior velocidade no
desenvolvimento, aproveitando o conhecimento prévio do autor na
linguagem. A versão \textit{desktop} é compreendida como uma etapa
intermediária no ciclo de vida do projeto, sendo a portabilidade para
Android e a integração de Realidade Aumentada objetivos para versões
futuras (Seção~\ref{sec:trabalhos-futuros}).

\textbf{Uso de inteligência artificial como ferramenta auxiliar.}
No desenvolvimento do código-fonte, foi utilizado o assistente de
inteligência artificial \textit{Claude} (Anthropic)~\cite{claude2024}
como ferramenta de suporte à programação. O uso restringiu-se à fase de
implementação, abrangendo: sugestões de estrutura modular, identificação
e correção de erros (\textit{debugging}), discussão de soluções para
compatibilidade multiplataforma, e revisão de boas práticas em Python.
Todo código gerado ou sugerido foi analisado, compreendido e validado
pelo autor antes de ser incorporado ao projeto. Conforme recomendações da
Academia Brasileira de Ciências~\cite{abc2023}, o histórico completo de
desenvolvimento está disponível no repositório público do projeto, com
todos os \textit{commits} documentados. Os prompts e trechos de
interação relevantes são apresentados no Apêndice~\ref{ap:ia}.

As etapas do desenvolvimento foram:
\begin{enumerate}
  \item Estudo e seleção de ferramentas e bibliotecas;
  \item Implementação incremental dos módulos do sistema;
  \item Testes funcionais de cada módulo de forma isolada;
  \item Integração dos módulos e testes de sistema;
  \item Avaliação com usuários (Capítulo~\ref{cap:resultados});
  \item Revisão, documentação e escrita da monografia.
\end{enumerate}

\section{Organização e Estrutura}

Este trabalho está organizado em cinco capítulos:

\begin{itemize}
  \item \textbf{Capítulo 1 --- Introdução:} apresenta as justificativas,
        objetivos, metodologia e estrutura do trabalho;
  \item \textbf{Capítulo 2 --- Referencial Teórico:} discute os
        fundamentos que embasam o trabalho, incluindo gamificação,
        educação musical, protocolo MIDI e trabalhos relacionados;
  \item \textbf{Capítulo 3 --- Desenvolvimento:} descreve a arquitetura
        do sistema, as decisões de projeto e a implementação dos módulos;
  \item \textbf{Capítulo 4 --- Resultados:} apresenta os resultados da
        avaliação com usuários;
  \item \textbf{Capítulo 5 --- Considerações Finais:} traz as conclusões,
        limitações e propostas de trabalhos futuros.
\end{itemize}

% ─────────────────────────────────────────────────────────────
%  CAPÍTULO 2 — REFERENCIAL TEÓRICO
% ─────────────────────────────────────────────────────────────
\chapter{Referencial Teórico}

\section{Gamificação no Ensino}

O termo \textit{gamificação} (\textit{gamification}) foi formalizado na
literatura acadêmica por \textcite{deterding2011}, que o definem como
``o uso de elementos de \textit{design} de jogos em contextos que não
são jogos''. A gamificação não consiste na criação de jogos completos,
mas na incorporação de mecanismos pontuais---como sistemas de pontos,
\textit{rankings}, conquistas, progressão em níveis e
\textit{feedback} imediato---a atividades cotidianas ou educacionais,
com o objetivo de aumentar o engajamento e a motivação dos participantes.

\textcite{kapp2012} argumenta que a gamificação é eficaz porque utiliza
elementos naturalmente motivadores presentes nos jogos---desafio
progressivo, recompensa, senso de progresso e autonomia---para manter o
aprendiz engajado em tarefas que, de outra forma, poderiam ser percebidas
como monótonas ou difíceis. \textcite{hamari2014}, em revisão sistemática
de 24 estudos empíricos sobre gamificação, concluíram que a maioria dos
estudos reporta efeitos positivos sobre motivação e engajamento, embora
ressaltem a importância do contexto e da qualidade do \textit{design} dos
elementos utilizados.

A prática de um instrumento musical, que exige repetição sistemática e
paciência, é um contexto especialmente favorável à aplicação da
gamificação. O \textit{feedback} visual imediato---saber se a nota tocada
foi correta e no momento certo---é um dos elementos gamificados mais
alinhados às demandas do aprendizado musical.

\section{Educação Musical e Aprendizado de Instrumentos}

O aprendizado de instrumentos musicais, em especial do piano, é
reconhecido como uma atividade multissensorial que envolve a integração
de habilidades motoras, auditivas e cognitivas. \textcite{hallam2015}
destaca que a prática regular de um instrumento está associada a
melhorias mensuráveis em funções executivas, memória de trabalho e
habilidades matemáticas. \textcite{miendlarzewska2014} revisaram
evidências neurológicas e comportamentais que indicam que o treinamento
musical modifica a estrutura e a função do cérebro, com efeitos mais
pronunciados quando iniciado na infância, mas presentes ao longo de
toda a vida.

Para o aprendizado efetivo, a literatura aponta dois fatores críticos:
a \textbf{qualidade do \textit{feedback}} e a
\textbf{frequência da prática}. O \textit{feedback} imediato e preciso
é fundamental para a consolidação do aprendizado motor~\cite{johansson2015}:
o aprendiz precisa saber, instante a instante, se sua execução está
correta para corrigir erros e reforçar acertos.

Nesse ponto, sistemas computacionais têm vantagem sobre o ensino
tradicional: podem fornecer \textit{feedback} instantâneo, preciso e
consistente, independente da disponibilidade de um professor. A
popularização de teclados MIDI USB de baixo custo reforça essa
oportunidade, tornando viável o desenvolvimento de sistemas educacionais
que complementem ou substituam, em parte, o ensino formal de piano.

\section{O Protocolo MIDI}

O protocolo \textit{MIDI} (\textit{Musical Instrument Digital Interface})
foi desenvolvido em 1983 como um padrão de comunicação digital entre
instrumentos musicais eletrônicos~\cite{midi1996}. Diferentemente de
formatos de áudio como MP3 ou WAV, o MIDI não armazena sons: armazena
\textbf{eventos musicais}---quando uma nota foi pressionada (\textit{note
on}), com qual intensidade (\textit{velocity}), e quando foi solta
(\textit{note off}).

\subsection{Estrutura de um arquivo MIDI}

Um arquivo MIDI padrão (\textit{Standard MIDI File} --- SMF) é composto
por:

\begin{itemize}
  \item \textbf{Cabeçalho (\textit{Header Chunk}):} define o formato do
        arquivo (0, 1 ou 2), o número de trilhas e a resolução temporal
        (em \textit{ticks} por pulso ou por quadro de tempo);
  \item \textbf{Trilhas (\textit{Track Chunks}):} sequências de eventos
        MIDI com seus respectivos \textit{delta times}---valores de tempo
        relativos ao evento anterior;
  \item \textbf{Mensagens de nota:}
        \begin{itemize}
          \item \texttt{note\_on} (canal, nota, velocidade): indica que
                uma tecla foi pressionada;
          \item \texttt{note\_off} (canal, nota): indica que a tecla foi
                solta (ou \texttt{note\_on} com velocidade zero);
          \item \texttt{control\_change}: transmite controles como pedal
                de sustain (CC\,64).
        \end{itemize}
\end{itemize}

\subsection{Numeração de notas MIDI}

As notas MIDI são codificadas como inteiros de 0 a 127. A nota 60
corresponde ao Dó central (C4); cada semitom corresponde a um incremento
de 1. Essa representação permite a conversão direta entre o valor MIDI
e a posição física de uma tecla em um piano padrão, como mostra a
equação~\eqref{eq:nota-oitava}:

\begin{equation}
  \text{oitava} = \left\lfloor \frac{n}{12} \right\rfloor - 1, \quad
  \text{semitom} = n \bmod 12
  \label{eq:nota-oitava}
\end{equation}

onde $n$ é o número da nota MIDI ($0 \leq n \leq 127$).

\subsection{MIDI no contexto educacional}

A ampla adoção do formato MIDI em \textit{softwares} de notação musical
gratuitos como o \textit{MuseScore}~\cite{musescore} tornou disponível
um vasto repertório de partituras digitais em formato SMF. Isso significa
que o usuário do sistema desenvolvido tem acesso imediato a milhares de
músicas, desde peças clássicas até arranjos populares, sem necessidade
de transcrição manual.

\section{Trabalhos Relacionados}

Diversos sistemas foram desenvolvidos para auxiliar o aprendizado de
piano com feedback visual e gamificação. A
Tabela~\ref{tab:relacionados} apresenta uma comparação entre os
principais trabalhos relacionados e o sistema proposto.

\begin{table}[htbp]
  \centering
  \caption{Comparação entre trabalhos relacionados e o sistema proposto.}
  \label{tab:relacionados}
  \resizebox{\textwidth}{!}{%
  \begin{tabular}{@{}lllccccc@{}}
    \toprule
    \textbf{Sistema} & \textbf{Plataforma} & \textbf{Gratuito} &
    \textbf{MIDI} & \textbf{Gamif.} & \textbf{Cód. aberto} &
    \textbf{Feedback visual} & \textbf{Feedback auditivo} \\
    \midrule
    Guitar Hero\textregistered & Console/PC & Não & Não & Sim & Não & Sim & Sim \\
    Synthesia       & PC/Mac      & Parcial & Sim & Parcial & Não & Sim & Sim \\
    PianoVision\textregistered & Meta Quest & Não & Sim & Parcial & Não & Sim (RA) & Sim \\
    Simply Piano    & iOS/Android & Parcial & Não & Sim & Não & Sim & Sim \\
    Yousician       & Multi       & Parcial & Não & Sim & Não & Sim & Sim \\
    \textbf{MIDI Hero} & \textbf{Linux/Win/Mac} & \textbf{Sim} &
    \textbf{Sim} & \textbf{Sim} & \textbf{Sim} & \textbf{Sim} & \textbf{Sim} \\
    \bottomrule
  \end{tabular}}
  \fonte{elaborado pelo autor}
\end{table}

\subsection{Guitar Hero\textregistered{} e Rock Band\textregistered}

\textit{Guitar Hero\textregistered} (Harmonix, 2005) é o precursor do
conceito de jogo musical com interface de notas descendentes
(\textit{falling notes}). O jogador utiliza um controle em formato de
guitarra e deve pressionar os botões corretos no momento em que as notas
atingem a zona de acerto na tela. Embora não utilize MIDI e não se
relacione diretamente com instrumentos reais, o modelo de interação
tornou-se referência para sistemas de prática musical gamificada.
Estudos apontaram que jogadores de \textit{Guitar Hero} frequentemente
desenvolvem interesse real em aprender a tocar guitarra após a
experiência com o jogo~\cite{hjorth2009}.

\subsection{Synthesia}

\textit{Synthesia}~\cite{synthesia} é um \textit{software} para PC e Mac
voltado especificamente ao aprendizado de piano. Lê arquivos MIDI e exibe
as notas como barras coloridas descendo em direção a um teclado virtual,
modelo muito similar ao adotado neste trabalho. Oferece versão gratuita
com funcionalidades limitadas e versão paga completa. Não é de código
aberto e não oferece sistema de pontuação gamificado.

\subsection{PianoVision\textregistered}

\textit{PianoVision\textregistered}~\cite{pianovision} é um aplicativo de
Realidade Aumentada para o \textit{headset} Meta Quest. Utiliza a câmera
do dispositivo para identificar as teclas de um piano físico e sobrepõe
animações de notas em RA, guiando o usuário visualmente. É o sistema
mais próximo da \textbf{visão original deste projeto}, porém requer
\textit{hardware} específico (Meta Quest, custo superior a
R\,\$\,3.000,00 no Brasil) e é pago, o que inviabiliza seu uso pela
população de menor renda.

\subsection{Simply Piano e Yousician}

\textit{Simply Piano}~\cite{simplypiano} e
\textit{Yousician}~\cite{yousician} são aplicativos para dispositivos
móveis que utilizam o microfone do \textit{smartphone} para reconhecer
as notas tocadas no piano físico, dispensando conexão MIDI. Oferecem
conteúdo gamificado com progressão em níveis e têm modelos
\textit{freemium}. Não são de código aberto.

\section{Ferramentas e Tecnologias Utilizadas}

\subsection{Python}

Python é uma linguagem de programação de alto nível, interpretada e de
propósito geral, reconhecida pela sintaxe clara e pela ampla
disponibilidade de bibliotecas científicas e de
multimídia~\cite{vanrossum2009}. A escolha do Python para este projeto
justifica-se pelo conhecimento prévio do autor na linguagem e pela
maturidade do ecossistema de bibliotecas disponíveis para processamento
de MIDI, síntese sonora e criação de interfaces gráficas.

\subsection{Pygame}

\textit{Pygame}~\cite{pygame2000} é uma biblioteca Python para
desenvolvimento de jogos e aplicações multimídia, construída sobre a
\textit{SDL} (\textit{Simple DirectMedia Layer}). Oferece recursos para
renderização de gráficos 2D, reprodução de áudio, captura de eventos e
gerenciamento de janelas. É multiplataforma (Linux, Windows, macOS) e de
código aberto.

\subsection{Mido e python-rtmidi}

\textit{Mido}~\cite{mido2013} é uma biblioteca Python para leitura,
escrita e envio de mensagens MIDI. Suporta o padrão SMF e a comunicação
com dispositivos MIDI físicos via \textit{python-rtmidi}, que provê o
\textit{backend} de acesso às portas MIDI do sistema operacional. É
utilizada neste projeto para o carregamento de arquivos MIDI e para a
leitura de dados provenientes do teclado MIDI do usuário.

\subsection{FluidSynth}

\textit{FluidSynth}~\cite{fluidsynth2000} é um sintetizador de áudio de
código aberto baseado no formato SoundFont (.sf2). É utilizado neste
projeto para duas finalidades: (i) a conversão de arquivos MIDI em áudio
(formato OGG) para reprodução sincronizada durante o jogo, e (ii) a
síntese sonora em tempo real em resposta às teclas pressionadas pelo
usuário, provendo retorno auditivo imediato.

\subsection{FFmpeg}

\textit{FFmpeg}~\cite{ffmpeg} é uma ferramenta de código aberto para
processamento de áudio e vídeo. É utilizada neste projeto na etapa de
conversão de arquivos WAV (gerados pelo FluidSynth) para o formato OGG,
que é então carregado pelo mixer do Pygame para reprodução durante o jogo.

% ─────────────────────────────────────────────────────────────
%  CAPÍTULO 3 — DESENVOLVIMENTO
% ─────────────────────────────────────────────────────────────
\chapter{Desenvolvimento}

\section{Visão Geral do Sistema}

O sistema desenvolvido é denominado \textbf{MIDI Hero} e consiste em
uma aplicação \textit{desktop} para prática de piano e teclado. O
usuário conecta um teclado MIDI ao computador via USB, seleciona uma
música no menu principal e deve pressionar as teclas corretas no momento
indicado pela interface visual gamificada. O sistema avalia a execução
em tempo real e atribui pontuação proporcional à precisão temporal e à
duração de cada nota sustentada.

O fluxo principal de execução é descrito a seguir:

\begin{enumerate}
  \item Na inicialização, o sistema verifica e converte automaticamente
        os arquivos MIDI disponíveis para o formato OGG;
  \item O usuário seleciona uma música no menu principal;
  \item O sistema carrega o arquivo MIDI e extrai os
        \textit{spans} de notas (pares início/fim de cada nota);
  \item Após contagem regressiva (3-2-1), o motor do jogo e a reprodução
        de áudio são iniciados de forma sincronizada;
  \item A interface gráfica renderiza barras descendentes representando
        as notas a tocar;
  \item Um \textit{thread} dedicado lê eventos do teclado MIDI em
        paralelo ao \textit{loop} de jogo;
  \item A cada nota pressionada, o sistema avalia a precisão temporal e
        inicia a pontuação progressiva;
  \item Ao término da música, uma tela de resultados exibe a pontuação
        final, a pontuação máxima possível e opções de recomeçar ou
        retornar ao menu.
\end{enumerate}

\section{Arquitetura Modular}

O sistema foi implementado como um pacote Python denominado
\texttt{midi\_game}, com oito módulos principais de responsabilidades
bem definidas. Essa separação facilita testes independentes de cada
componente e viabiliza extensões futuras sem impacto nos demais módulos.
A Tabela~\ref{tab:modulos} resume os módulos e suas responsabilidades.

\begin{table}[htbp]
  \centering
  \caption{Módulos do pacote \texttt{midi\_game} e suas responsabilidades.}
  \label{tab:modulos}
  \begin{tabular}{@{}ll@{}}
    \toprule
    \textbf{Módulo} & \textbf{Responsabilidade} \\
    \midrule
    \texttt{config.py}          & Gerenciamento de configurações (JSON) \\
    \texttt{midi\_loader.py}    & Leitura e \textit{parsing} de arquivos MIDI \\
    \texttt{game\_engine.py}    & Controle temporal do jogo \\
    \texttt{graphics.py}        & Interface gráfica (Pygame) \\
    \texttt{sounds.py}          & Reprodução de áudio sincronizado \\
    \texttt{midi\_input.py}     & Detecção e abertura de portas MIDI \\
    \texttt{live\_monitor.py}   & Síntese sonora em tempo real \\
    \texttt{audio\_converter.py}& Conversão MIDI $\to$ WAV $\to$ OGG \\
    \bottomrule
  \end{tabular}
  \fonte{elaborado pelo autor}
\end{table}

\subsection{Módulo de Configuração (\texttt{config.py})}

O módulo \texttt{config.py} gerencia as preferências do usuário,
armazenadas em um arquivo JSON (\texttt{config.json}) no mesmo diretório
do pacote. As configurações disponíveis são: número de teclas do teclado
(49, 61, 76 ou 88), nota MIDI mais grave, porta MIDI preferida,
habilitação do som ao vivo e volumes de música e teclado.

Um problema crítico identificado nas versões iniciais foi a definição
do caminho do arquivo de configuração como string relativa
(\texttt{"midi\_game/config.json"}), o que causava erros quando o sistema
era executado a partir de um diretório diferente da raiz do projeto. A
solução adotada utiliza a biblioteca \texttt{pathlib} em conjunto com a
variável especial \texttt{\_\_file\_\_} do Python para construir um
caminho absoluto baseado na localização do próprio módulo:

\begin{lstlisting}[caption={Definição de caminho absoluto com pathlib.},
                   label={lst:config-path}]
from pathlib import Path
CONFIG_PATH = Path(__file__).parent / "config.json"
\end{lstlisting}

Esse padrão foi replicado em todos os módulos que referenciam arquivos
em disco, garantindo portabilidade entre sistemas operacionais e
independência do diretório de execução.

\subsection{Módulo de Carregamento MIDI (\texttt{midi\_loader.py})}

O módulo \texttt{midi\_loader.py} lê e interpreta arquivos MIDI usando
a biblioteca \textit{Mido}. A função \texttt{load\_midi\_file()} lê o
arquivo e converte os \textit{delta times} relativos em tempos
acumulados em segundos, iterando sobre os eventos de todas as trilhas.

A função \texttt{build\_note\_spans()} converte a lista de eventos em
\textit{spans} de notas---dicionários contendo o número da nota MIDI, o
instante de início e o instante de fim. A função suporta sobreposição
da mesma nota (quando uma nota é pressionada novamente antes de ser
solta), utilizando uma pilha por nota:

\begin{lstlisting}[caption={Trecho da função build\_note\_spans().},
                   label={lst:spans}]
stacks: dict[int, list[float]] = {}  # note -> [start_times]
for e in events:
    if e["type"] == "note_on" and e["velocity"] > 0:
        stacks.setdefault(e["note"], []).append(e["time"])
    elif e["type"] == "note_off" or e["velocity"] == 0:
        stk = stacks.get(e["note"])
        if stk:
            start = stk.pop()
            if e["time"] > start:
                spans.append({"note": e["note"],
                              "start": start,
                              "end": e["time"]})
\end{lstlisting}

\subsection{Motor do Jogo (\texttt{game\_engine.py})}

A classe \texttt{GameEngine} gerencia o tempo de jogo. Ao ser iniciada
com o método \texttt{start()}, registra o \textit{timestamp} absoluto
de início. O método \texttt{get\_current\_time()} retorna o tempo atual
do jogo em segundos, com valores negativos durante o \textit{pré-roll}
(período anterior ao início da música, de 2 segundos), que permite ao
usuário visualizar as primeiras notas com antecedência.

O método \texttt{get\_visible\_spans()} filtra os spans que devem ser
renderizados em um dado instante, com uma janela de 8 segundos à frente
e 1 segundo atrás do instante atual do jogo, limitando o processamento
gráfico às notas relevantes para o momento.

\subsection{Interface Gráfica (\texttt{graphics.py})}

O módulo \texttt{graphics.py} implementa toda a interface visual do
jogo com Pygame em modo \textit{fullscreen}. Os elementos principais são:

\begin{itemize}
  \item \textbf{Teclado virtual:} representação gráfica do teclado de
        piano na parte inferior da tela, com teclas brancas e pretas.
        O número de teclas e a nota mais grave são configuráveis pelo
        usuário;
  \item \textbf{Barras de nota (\textit{note bars}):} retângulos
        coloridos que descem do topo da tela em direção ao teclado,
        representando as notas a pressionar. A largura corresponde à
        tecla do teclado; a altura é proporcional à duração da nota;
  \item \textbf{HUD (\textit{Heads-Up Display}):} exibe a pontuação
        atual e a pontuação máxima possível no topo da tela;
  \item \textbf{\textit{Popups} de pontuação:} exibem os pontos
        ganhos ao finalizar cada nota, com animação de flutuação e
        desvanecimento gradual.
\end{itemize}

\subsection{Entrada MIDI (\texttt{midi\_input.py})}

O módulo \texttt{midi\_input.py} gerencia a detecção e abertura de
portas MIDI. A função \texttt{pick\_best\_port()} implementa uma
heurística para selecionar automaticamente o dispositivo mais adequado:
prioriza a porta configurada pelo usuário; na sua ausência, seleciona
a primeira porta ``real'' disponível, excluindo portas virtuais do
sistema (como \textit{MIDI Through}) que não correspondem a instrumentos
físicos.

\subsection{Síntese Sonora em Tempo Real (\texttt{live\_monitor.py})}

O módulo \texttt{live\_monitor.py} implementa o retorno auditivo ao
vivo: quando o usuário pressiona uma tecla no teclado MIDI, o sistema
emite imediatamente o som correspondente via \textit{FluidSynth} e
SoundFont, tornando a experiência mais semelhante à de tocar um piano
real.

Um desafio de portabilidade foi a inicialização do driver de áudio do
FluidSynth. A solução detecta o sistema operacional em tempo de execução
e tenta os drivers na ordem adequada a cada plataforma:

\begin{lstlisting}[caption={Seleção de driver de áudio por plataforma.},
                   label={lst:drivers}]
_SYS = platform.system()
if _SYS == "Windows":
    DRIVERS = ("dsound", "portaudio", "sdl2")
elif _SYS == "Darwin":   # macOS
    DRIVERS = ("coreaudio", "portaudio", "sdl2")
else:                    # Linux
    DRIVERS = ("pulseaudio", "pipewire", "alsa", "jack", "portaudio", "sdl2")
\end{lstlisting}

\subsection{Conversor de Áudio (\texttt{audio\_converter.py})}

O módulo \texttt{audio\_converter.py} converte automaticamente os
arquivos MIDI para o formato OGG na primeira execução do sistema, em
duas etapas: (i) renderização MIDI $\to$ WAV usando o \textit{FluidSynth}
(via linha de comando); (ii) conversão WAV $\to$ OGG usando o
\textit{FFmpeg} (via biblioteca \textit{ffmpeg-python}). Se os arquivos
OGG já existirem, a conversão é ignorada, evitando reprocessamento
desnecessário.

O módulo inclui localização automática de SoundFonts em caminhos padrão
de cada sistema operacional e exibe instruções de instalação das
ferramentas necessárias caso não sejam encontradas, com comandos
específicos para cada plataforma (pacman, apt-get, brew, winget).

\section{Sistema de Pontuação}
\label{sec:pontuacao}

O sistema de pontuação foi projetado para recompensar tanto a
\textbf{precisão temporal} (acertar a nota no momento certo) quanto a
\textbf{persistência} (manter a tecla pressionada pelo tempo correto da
nota na partitura).

\subsection{Janela de tolerância e multiplicador de qualidade}

Quando o usuário pressiona uma tecla, o sistema verifica se existe uma
nota esperada dentro de uma janela de tolerância temporal de 80\,ms
(\texttt{HIT\_WINDOW}). Caso positivo, o \textit{span} é
``engatado'' e um multiplicador de qualidade $q$ é calculado com base
no erro temporal $\Delta t$, conforme a
Tabela~\ref{tab:multiplicadores}:

\begin{table}[htbp]
  \centering
  \caption{Multiplicadores de qualidade por precisão temporal.}
  \label{tab:multiplicadores}
  \begin{tabular}{@{}ccc@{}}
    \toprule
    \textbf{Erro temporal} & \textbf{Classificação} & \textbf{Multiplicador $q$} \\
    \midrule
    $\Delta t < 30\,\text{ms}$          & Perfeito & 1{,}0 \\
    $30 \leq \Delta t < 60\,\text{ms}$  & Bom      & 0{,}7 \\
    $60 \leq \Delta t \leq 80\,\text{ms}$ & Ok     & 0{,}4 \\
    \bottomrule
  \end{tabular}
  \fonte{elaborado pelo autor}
\end{table}

\subsection{Pontuação progressiva}

Os pontos são acumulados continuamente enquanto o usuário mantém a tecla
pressionada, proporcionalmente ao tempo em que a nota está sendo
corretamente executada. A taxa de pontuação de cada span é calculada por:

\begin{equation}
  r = \frac{P_{\max}}{d_{\text{span}}}
  \label{eq:taxa}
\end{equation}

onde $P_{\max}$ é a pontuação máxima por span (100 pontos) e
$d_{\text{span}}$ é a duração do span em segundos. A cada
\textit{frame}, os pontos acrescidos são:

\begin{equation}
  \Delta p = r \cdot \Delta t_{\text{frame}} \cdot q
  \label{eq:delta-pontos}
\end{equation}

onde $\Delta t_{\text{frame}}$ é o intervalo de tempo entre frames e $q$
é o multiplicador de qualidade definido no engajamento. O total acumulado
é limitado a $P_{\max} \cdot q$ por span, impedindo que o usuário que
pressiona a tecla com atraso, mas a mantém por tempo excedente, supere a
pontuação máxima.

\section{Compatibilidade Multiplataforma}

Um dos objetivos específicos do trabalho foi garantir funcionamento
correto nos três principais sistemas operacionais. Os problemas de
portabilidade identificados e as soluções adotadas são descritos na
Tabela~\ref{tab:portabilidade}:

\begin{table}[htbp]
  \centering
  \caption{Problemas de portabilidade e soluções adotadas.}
  \label{tab:portabilidade}
  \resizebox{\textwidth}{!}{%
  \begin{tabular}{@{}p{4.5cm}p{4.5cm}p{5cm}@{}}
    \toprule
    \textbf{Problema} & \textbf{Causa} & \textbf{Solução} \\
    \midrule
    Caminhos de arquivo inválidos fora do diretório raiz &
    Caminhos relativos hardcoded &
    \texttt{pathlib.Path(\_\_file\_\_).parent} em todos os módulos \\
    \midrule
    Falha na inicialização do áudio (FluidSynth) &
    Drivers diferentes por SO &
    Lista de drivers priorizada por plataforma em runtime \\
    \midrule
    SoundFont não encontrado &
    Caminhos de sistema diferentes por SO &
    Busca em caminhos padrão de Linux, macOS e Windows \\
    \midrule
    Erro de sintaxe em Python < 3.10 &
    Uso de \texttt{X | Y} em anotações de tipo &
    \texttt{from \_\_future\_\_ import annotations} em todos os módulos \\
    \bottomrule
  \end{tabular}}
  \fonte{elaborado pelo autor}
\end{table}

\section{Instalação e Distribuição}

Para facilitar a instalação em diferentes plataformas, foram criados
\textit{scripts} específicos:

\begin{itemize}
  \item \textbf{Makefile} (Linux e macOS): cria ambiente virtual Python
        (\textit{venv}), instala as dependências listadas no
        \texttt{requirements.txt} e exibe instruções para instalação das
        ferramentas de sistema. O jogo é iniciado com o comando
        \texttt{make run};
  \item \textbf{install.bat / run.bat} (Windows): equivalentes ao
        Makefile para o ambiente Windows, com instruções específicas para
        instalação via \texttt{winget} do FluidSynth e do FFmpeg.
\end{itemize}

As dependências Python do projeto estão declaradas no arquivo
\texttt{requirements.txt} com versões mínimas:

\begin{lstlisting}[caption={requirements.txt do projeto MIDI Hero.},
                   label={lst:requirements},language={}]
pygame>=2.5.0
pyfluidsynth>=1.3.3
mido>=1.3.0
python-rtmidi>=1.5.0
ffmpeg-python>=0.2.0
\end{lstlisting}

\section{Validação Automatizada --- CI/CD}

Para garantir a qualidade do código em alterações futuras, foi
configurado um \textit{pipeline} de integração contínua (CI) utilizando
o \textit{GitHub Actions}. O \textit{pipeline} é composto por três
etapas executadas automaticamente em todo \textit{push} ou
\textit{pull request}:

\begin{enumerate}
  \item \textbf{Verificação de sintaxe} (\texttt{py\_compile}): executa
        nos três sistemas operacionais (Ubuntu, Windows, macOS) com
        Python 3.10, 3.11 e 3.12, totalizando 9 combinações;
  \item \textbf{Instalação e importação} (Linux): instala as
        dependências de sistema via \texttt{apt-get}, instala o
        \texttt{requirements.txt}, verifica que todos os caminhos são
        absolutos e que todos os módulos importam corretamente;
  \item \textbf{Verificação do \texttt{pip install}} (Windows e macOS):
        valida que o \texttt{requirements.txt} instala sem erros nessas
        plataformas.
\end{enumerate}

% ─────────────────────────────────────────────────────────────
%  CAPÍTULO 4 — RESULTADOS
% ─────────────────────────────────────────────────────────────
\chapter{Resultados}
\label{cap:resultados}

% ── PLACEHOLDER ──────────────────────────────────────────────
% Esta seção deve ser preenchida após a realização dos testes
% com usuários. Os títulos, questionário SUS e estrutura
% sugerida estão prontos; substitua os blocos [PREENCHER]
% pelos dados reais coletados.
% ─────────────────────────────────────────────────────────────

\section{Protocolo de Avaliação}

A avaliação do sistema foi conduzida por meio de um experimento com
usuários voluntários sem experiência prévia com piano ou teclado
musical. O protocolo adotado compreendeu as seguintes etapas:

\begin{enumerate}
  \item \textbf{Familiarização (10 min):} o participante explorou
        livremente o sistema, navegou pelos menus e realizou uma
        execução sem avaliação formal;
  \item \textbf{Sessão de jogo (5 min):} o participante executou a
        música \textit{[PREENCHER: nome da música utilizada]} no nível
        padrão do sistema;
  \item \textbf{Questionário pós-uso:} aplicado imediatamente após a
        sessão, composto por (i) a escala \textit{System Usability
        Scale} (SUS)~\cite{brooke1996}, com 10 itens em escala Likert
        de 5 pontos, e (ii) três perguntas abertas sobre experiência
        anterior com música, percepção de aprendizado e sugestões de
        melhoria.
\end{enumerate}

A escala SUS produz uma pontuação entre 0 e 100 pontos, em que valores
acima de 68 são considerados acima da média de usabilidade e valores
acima de 80 são classificados como ``bons'' ou ``excelentes''
\cite{brooke1996}.

\section{Perfil dos Participantes}

% ── PREENCHER ────────────────────────────────────────────────
Participaram do estudo \textbf{[N]} voluntários, recrutados
\textbf{[PREENCHER: descrever forma de recrutamento]}.
A Tabela~\ref{tab:participantes} resume o perfil demográfico da amostra.

\begin{table}[htbp]
  \centering
  \caption{Perfil dos participantes da avaliação.}
  \label{tab:participantes}
  \begin{tabular}{@{}lc@{}}
    \toprule
    \textbf{Característica} & \textbf{Valor} \\
    \midrule
    Número de participantes          & [N]  \\
    Faixa etária                     & [PREENCHER] \\
    Gênero (M/F/Outro)               & [PREENCHER] \\
    Sem experiência musical formal   & [N]\% \\
    Com experiência musical (não piano) & [N]\% \\
    Familiaridade com jogos digitais & [PREENCHER] \\
    \bottomrule
  \end{tabular}
  \fonte{elaborado pelo autor}
\end{table}

\section{Resultados Quantitativos}

\subsection{Pontuação no jogo}

% ── PREENCHER ────────────────────────────────────────────────
A Tabela~\ref{tab:pontuacoes} apresenta as pontuações obtidas pelos
participantes, expressas como percentual da pontuação máxima possível
da música selecionada.

\begin{table}[htbp]
  \centering
  \caption{Pontuações obtidas pelos participantes (percentual da
           pontuação máxima).}
  \label{tab:pontuacoes}
  \begin{tabular}{@{}lccc@{}}
    \toprule
    \textbf{Estatística} & \textbf{Valor (\%)} \\
    \midrule
    Média    & [PREENCHER] \\
    Mediana  & [PREENCHER] \\
    Desvio padrão & [PREENCHER] \\
    Mínimo   & [PREENCHER] \\
    Máximo   & [PREENCHER] \\
    \bottomrule
  \end{tabular}
  \fonte{elaborado pelo autor}
\end{table}

\subsection{Usabilidade (SUS)}

% ── PREENCHER ────────────────────────────────────────────────
A pontuação média na escala SUS obtida foi de
\textbf{[PREENCHER]} pontos (desvio padrão: [PREENCHER]),
o que corresponde à classificação \textbf{[PREENCHER]}
segundo os critérios de \textcite{brooke1996}.

% Quando tiver os dados, substitua o parágrafo acima por algo como:
% ``A pontuação média na escala SUS foi de \textbf{82,5} pontos
%   (dp = 7,3), classificada como ``boa'' segundo Brooke (1996),
%   indicando que o sistema apresenta boa usabilidade percebida
%   pelos participantes.''

\section{Resultados Qualitativos}

% ── PREENCHER ────────────────────────────────────────────────
As respostas às perguntas abertas do questionário foram analisadas por
análise de conteúdo. Os principais temas identificados foram:

\begin{itemize}
  \item \textbf{Percepção de aprendizado:} [PREENCHER --- ex.: ``a
        maioria dos participantes relatou perceber melhora na
        coordenação mão-tela após a sessão''];
  \item \textbf{Pontos positivos:} [PREENCHER --- ex.: feedback visual
        claro, sistema de pontuação motivador];
  \item \textbf{Pontos de melhoria:} [PREENCHER --- ex.: ausência de
        modo para iniciantes com velocidade reduzida, necessidade de
        tutorial].
\end{itemize}

\section{Discussão}

% ── PREENCHER ────────────────────────────────────────────────
Os resultados indicam que [PREENCHER]. A pontuação SUS de [PREENCHER]
sugere que [PREENCHER]. As limitações observadas---em especial
[PREENCHER]---apontam para melhorias prioritárias descritas nas
propostas de trabalhos futuros (Seção~\ref{sec:trabalhos-futuros}).

% ─────────────────────────────────────────────────────────────
%  CAPÍTULO 5 — CONSIDERAÇÕES FINAIS
% ─────────────────────────────────────────────────────────────
\chapter{Considerações Finais}

\section{Conclusões}

Este trabalho apresentou o desenvolvimento do \textbf{MIDI Hero}, um
sistema computacional gamificado para suporte à prática de piano e
teclado. O sistema é capaz de ler partituras no formato MIDI, gerar
feedback visual sincronizado em tempo real por meio de barras de notas
descendentes, avaliar a execução do usuário por sistema de pontuação
progressiva e prover retorno auditivo imediato via síntese sonora por
SoundFont.

Os objetivos específicos propostos foram atendidos:

\begin{itemize}
  \item A leitura de arquivos MIDI foi implementada com suporte a
        sobreposição de notas e conversão de \textit{delta times} para
        tempo absoluto em segundos;
  \item A interface gráfica gamificada foi desenvolvida com teclado
        virtual, barras de notas, HUD e animações de pontuação;
  \item O sistema de pontuação progressiva foi implementado,
        recompensando precisão temporal e duração correta das notas;
  \item A síntese sonora em tempo real foi integrada via FluidSynth,
        com detecção automática de SoundFonts em múltiplos caminhos;
  \item A compatibilidade multiplataforma foi garantida para Linux,
        Windows e macOS, com instaladores dedicados para cada plataforma;
  \item A arquitetura modular adotada cria uma base extensível para
        versões futuras do sistema.
\end{itemize}

Do ponto de vista da engenharia de \textit{software}, a experiência
reforçou a importância de tratar portabilidade como requisito de
primeira classe---e não como preocupação posterior---no desenvolvimento
de sistemas desktop. Problemas simples como caminhos relativos e ordens
de driver de áudio, invisíveis durante o desenvolvimento em um único
sistema operacional, tornam-se críticos na distribuição multiplataforma.

O uso de inteligência artificial como ferramenta auxiliar no
desenvolvimento mostrou-se produtivo, especialmente na identificação de
problemas de portabilidade e na discussão de boas práticas de
arquitetura. A experiência corrobora a perspectiva de
\textcite{kapp2012}: ferramentas que automatizam tarefas operacionais
deslocam o foco do desenvolvedor para decisões de nível mais alto,
aumentando a qualidade do resultado final.

\section{Limitações}

As principais limitações do sistema em sua versão atual são:

\begin{itemize}
  \item \textbf{Dependência de hardware MIDI:} o sistema requer um
        teclado MIDI para captura da execução. Não há suporte a
        reconhecimento de áudio por microfone, o que limita o acesso
        a usuários sem esse equipamento;
  \item \textbf{Plataforma exclusivamente \textit{desktop}:} a versão
        atual não está disponível para dispositivos móveis, que
        representam o principal vetor de acesso à tecnologia para
        populações de menor renda;
  \item \textbf{Dependências de sistema:} a conversão MIDI $\to$ OGG
        requer FluidSynth e FFmpeg instalados no sistema, o que
        pode ser um obstáculo para usuários sem experiência técnica;
  \item \textbf{Ausência de currículo progressivo:} o sistema não
        orienta o usuário sobre qual música praticar primeiro,
        deixando a seleção de repertório inteiramente a cargo do
        usuário.
\end{itemize}

\section{Trabalhos Futuros}
\label{sec:trabalhos-futuros}

A arquitetura modular adotada foi projetada com as seguintes extensões
em vista:

\begin{itemize}
  \item \textbf{Portabilidade para Android:} o desenvolvimento de uma
        versão para dispositivos móveis Android, com interface adaptada
        a telas sensíveis ao toque, expandiria significativamente o
        alcance do sistema;

  \item \textbf{Realidade Aumentada:} a integração de RA---utilizando
        a câmera do dispositivo para sobrepor animações de notas sobre
        um teclado físico real---representa o próximo passo natural de
        evolução do projeto, retomando a visão original. Sistemas como
        o PianoVision\textregistered~\cite{pianovision} demonstram a
        viabilidade técnica dessa abordagem;

  \item \textbf{Reconhecimento de áudio:} a adição de detecção de
        notas por microfone eliminaria a dependência de hardware MIDI,
        democratizando o acesso ao sistema;

  \item \textbf{Currículo progressivo e sistema de níveis:} a
        estruturação de músicas em níveis de dificuldade com progressão
        guiada transformaria o sistema em uma plataforma educacional
        completa;

  \item \textbf{Modo multiplayer e placar online:} funcionalidades de
        competição e colaboração poderiam aumentar o engajamento e a
        retenção a longo prazo.
\end{itemize}

% ──────────────────────────────────────────────────────────────
%  PÓS-TEXTUAIS
% ──────────────────────────────────────────────────────────────
\postextual

\printbibliography[title=Referências]

% ─────────────────────────────────────────────────────────────
%  APÊNDICES
% ─────────────────────────────────────────────────────────────
\begin{apendicesenv}
\partapendices

% ── Apêndice A — Questionário SUS ────────────────────────────
\chapter{Questionário de Avaliação --- Escala SUS}
\label{ap:sus}

O questionário abaixo foi aplicado aos participantes imediatamente
após a sessão de uso do sistema. A escala SUS (\textit{System Usability
Scale})~\cite{brooke1996} é composta por 10 afirmações respondidas em
escala Likert de 5 pontos (1 = Discordo totalmente; 5 = Concordo
totalmente). Os itens ímpares têm sentido positivo; os pares, negativo.

\begin{enumerate}
  \item Eu gostaria de usar este sistema com frequência.
  \item Achei o sistema desnecessariamente complexo.
  \item Achei o sistema fácil de usar.
  \item Acho que precisaria de ajuda de uma pessoa técnica para usar
        este sistema.
  \item Achei as várias funções do sistema bem integradas.
  \item Achei que havia muita inconsistência no sistema.
  \item Imagino que a maioria das pessoas aprenderia a usar este
        sistema rapidamente.
  \item Achei o sistema muito difícil de usar.
  \item Me senti muito confiante usando o sistema.
  \item Precisei aprender muitas coisas antes de conseguir usar
        este sistema.
\end{enumerate}

\noindent
\textbf{Perguntas abertas:}
\begin{enumerate}
  \item[11.] Você tem experiência com algum instrumento musical?
        Qual?
  \item[12.] Após usar o MIDI Hero, você sente que aprendeu algo
        sobre o instrumento ou sobre leitura de notas? Descreva
        brevemente.
  \item[13.] O que você mudaria ou adicionaria ao sistema?
\end{enumerate}

% ── Apêndice B — Uso de IA no desenvolvimento ────────────────
\chapter{Uso de Inteligência Artificial no Desenvolvimento}
\label{ap:ia}

Em conformidade com as recomendações da Academia Brasileira de
Ciências~\cite{abc2023} e com o princípio de transparência acadêmica,
este apêndice documenta o uso do assistente de inteligência artificial
\textit{Claude} (Anthropic)~\cite{claude2024} como ferramenta auxiliar
no desenvolvimento do código-fonte deste trabalho.

\section*{Escopo do uso}

O uso da ferramenta restringiu-se integralmente à fase de
implementação do \textit{software}, nas seguintes atividades:

\begin{itemize}
  \item Discussão de alternativas para a estrutura modular do projeto;
  \item Identificação e correção de erros de portabilidade
        (caminhos relativos, ordem de drivers de áudio);
  \item Sugestão de uso de \texttt{pathlib} para caminhos absolutos;
  \item Identificação de incompatibilidade de sintaxe de anotações de
        tipo com Python < 3.10 e solução via
        \texttt{from \_\_future\_\_ import annotations};
  \item Revisão e discussão de boas práticas em Python.
\end{itemize}

\section*{O que \emph{não} foi feito por IA}

\begin{itemize}
  \item O texto dissertativo desta monografia não foi gerado por IA;
  \item A concepção do sistema, a definição de requisitos, a escolha
        de tecnologias e as decisões de arquitetura foram tomadas
        pelo autor;
  \item Todo código gerado ou sugerido pela ferramenta foi
        analisado, compreendido e validado pelo autor antes de
        ser incorporado ao projeto.
\end{itemize}

\section*{Rastreabilidade}

O histórico completo de desenvolvimento---incluindo todos os
\textit{commits} com mensagens descritivas---está disponível no
repositório público do projeto no GitHub:
\url{https://github.com/samuelsilvaop/tcc}.
Os \textit{commits} originados de interações com a ferramenta de IA
contêm em sua mensagem a referência à sessão correspondente.

\end{apendicesenv}

% ─────────────────────────────────────────────────────────────
%  ÍNDICE REMISSIVO
% ─────────────────────────────────────────────────────────────
\printindex

\end{document}
